% =============================================================
%  Test Document: Combined GPU/Corpus Documentation with TikZ Reference
%  Standalone document with custom styling
% =============================================================
\documentclass[11pt,a4paper]{article}

% Basic packages
\usepackage{listings}
\usepackage{tikz}
\usetikzlibrary{arrows.meta,positioning,calc,shapes.geometric,fit,decorations.pathreplacing}
\usepackage{circuitikz}
\usepackage{xcolor}
\usepackage{geometry}
\usepackage{hyperref}

% Define brand colors
\definecolor{atLime}{HTML}{CEEB11}
\definecolor{atBlack}{HTML}{121212}
\definecolor{atInk}{HTML}{232323}
\definecolor{atPaper}{HTML}{FFFFFF}
\definecolor{atGray}{HTML}{5C5C5C}
\definecolor{atPanel}{HTML}{272B2B}
\definecolor{atSand}{HTML}{E3E2DA}
\definecolor{atStone}{HTML}{C9C7C2}
\definecolor{atBorderStrong}{HTML}{C9C7C2}
\definecolor{atBg}{HTML}{F7F7F6}
\definecolor{atBgSubtle}{HTML}{F2F2F2}
\definecolor{atBgMuted}{HTML}{E3E2DA}
\definecolor{atAccent}{HTML}{CEEB11}
\definecolor{atFg}{HTML}{232323}
\definecolor{atFgMuted}{HTML}{5C5C5C}

\geometry{margin=1in}

% ---- Shell / code listing style ------------------------------
\lstdefinestyle{atshell}{
  basicstyle=\ttfamily\small\color{atFg},
  backgroundcolor=\color{atBgSubtle},
  frame=single,
  rulecolor=\color{atBorderStrong},
  framesep=5pt,
  breaklines=true,
  breakindent=0pt,
  postbreak=\mbox{\textcolor{atBorderStrong}{$\hookrightarrow$}\space},
  columns=flexible,
  keepspaces=true,
  showstringspaces=false,
  commentstyle=\color{atFgMuted}\itshape,
  morecomment=[l]{\#},
  xleftmargin=0pt,
  aboveskip=8pt,
  belowskip=8pt,
}
\lstset{style=atshell}

\newcommand{\cmd}[1]{\texttt{\small #1}}

% ---- Document metadata ----------------------------------------
\title{Test Document: GPU Pool and Corpus System Documentation}
\subtitle{Combined reference: LLM deployment, document retrieval, and transmission-line TikZ diagram}
\author{Eugen Starschich}
\date{\today}

\begin{document}

\maketitle

\begin{abstract}
End-to-end account of how a local large language model is deployed on the shared
GPU pool --- how GPUs are borrowed, why the work happens inside a container, where
files must live to survive, how the model runtime is installed into an image
without package manager access, and how the resulting endpoint is reached from a
remote machine --- followed by \textit{corpus}, the retrieval
system that answers questions from our own documents over that endpoint.
Plus: TikZ representation of transmission line model.
\end{abstract}

\tableofcontents
\newpage

% ============================================================
%  PART I: GPU POOL DOCUMENTATION
% ============================================================

\part{GPU Pool and LLM Deployment}

\section{Architecture}
\label{sec:architecture}

Two machines are involved: the developer VM (dev-eugen) and the GPU host
(camimo). The pool claim lifecycle is explicit (claim $\rightarrow$ release), the
container home is ephemeral (data must live on /shared), and the tunnel
forwarding the endpoint must survive container recreation.

\subsection{The Shared GPU Pool}

The GPU pool is not a Kubernetes cluster. It is a single physical box
(camimo) with four GPUs, plus a DRA (Dynamic Resource Allocation) script
that hands one GPU at a time to developers on first-come-first-served basis. The claim lifecycle is explicit:

\begin{lstlisting}
# --- pool -----------------------------------------------------------
gpu status              # see current claim holder
gpu claim 1             # claim GPU 1 (blocks forever until release)
gpu shell               # interactive shell inside the container
gpu shell /path/to/script.sh    # run command and attach
gpu release             # release GPU; container destroyed
\end{lstlisting}

Claims do not expire, by design. A forgotten claim blocks the pool indefinitely.

\section{Storage Model}

Only /shared spans both machines. The container home is the GPU host's, not the
developer VM's.

\begin{table}[!ht]
  \centering
  \renewcommand{\arraystretch}{1.25}
  \begin{tabularx}{\linewidth}{@{}l X@{}}
    \toprule
    \textbf{Location} & \textbf{Purpose} \\\midrule
    \cmd{/shared/} & Persistent data (model weights, scripts, logs) \\\midrule
    \cmd{container home} & Ephemeral (deleted on release) \\\midrule
    \cmd{~/work/raw} & Dashboard uploads (dev-eugen only) \\\bottomrule
  \end{tabularx}
  \caption{Storage model: what survives and where.}
  \label{tab:storage}
\end{table}

\section{Ollama Runtime}

Install once, then the model and the container survive independently.

\begin{lstlisting}
# --- inside the container -------------------------------------------
nvidia-smi                            # verify GPU is claimed
ps -eo pid,etime,cmd | grep "[b]in/ollama serve"
tail -f /shared/$USER/ollama-serve.log
\end{lstlisting}

The script /shared/$USER/start-ollama.sh survives release. After a later claim:

\begin{lstlisting}
gpu claim 1
gpu shell /shared/$USER/start-ollama.sh
\end{lstlisting}

\section{Model and Endpoint}

\begin{lstlisting}
# --- model ----------------------------------------------------------
ollama list
ollama pull  gpt-oss:120b
ollama run   gpt-oss:120b --verbose "prompt"
ollama stop  gpt-oss:120b       # unload from VRAM, keep on disk

# --- endpoint -------------------------------------------------------
curl http://localhost:11435/api/tags
curl http://localhost:11435/api/generate \
  -d '{"model":"gpt-oss:120b","prompt":"...","stream":false}'
\end{lstlisting}

The local port (11435) differs from the default (11434) because the relay re-quotes
arguments; nested quoting does not survive.

\section{Remote Access (Laptop $\rightarrow$ GPU Host)}

From the developer VM, reach the GPU host over the office LAN. The SSH tunnel
forwards the model endpoint so the laptop can talk to it as if it were local.

\begin{lstlisting}
ssh -f -N -L 127.0.0.1:11434:172.17.0.4:11434 camimo
\end{lstlisting}

Confirm the bridge address first, as it may differ between claims:

\begin{lstlisting}
ssh camimo "docker inspect -f \
  '{{range .NetworkSettings.Networks}}{{.IPAddress}}{{end}}' gpu-eugen"
\end{lstlisting}

\section{Pitfalls}

\begin{table}[!ht]
  \centering
  \renewcommand{\arraystretch}{1.25}
  \begin{tabularx}{\linewidth}{@{}>{\raggedright\arraybackslash}p{4.4cm} X@{}}
    \toprule
    \textbf{Symptom} & \textbf{Cause and remedy} \\\midrule
    Claim never released &
    Claims do not expire, by design. A forgotten claim blocks the pool
    indefinitely. Release explicitly. \\[2pt]

    \cmd{Text file busy} on a freshly staged binary &
    NFS write-back lag on a new executable under \cmd{/shared}. It clears by
    itself; wait rather than re-copying. \\[2pt]

    \cmd{pgrep -f "ollama serve"} always matches &
    The search pattern appears in the searching command's own command line.
    Use \cmd{grep "[b]in/ollama serve"}. \\[2pt]

    Quoted one-liner silently does nothing &
    The relay re-quotes arguments; nested quoting does not survive. Put the
    commands in a script on \cmd{/shared}. \\[2pt]

    \cmd{curl} returns the wrong model &
    A local Ollama already occupies port 11434, so the forward failed to bind
    and the query answered locally. Forward to a free local port such as
    11435. \\[2pt]

    \cmd{ssh -N} appears to hang &
    Expected. \cmd{-N} runs no remote command; it holds the tunnel open and
    prints nothing. Test from a second terminal. \\[2pt]

    Work lost on release &
    Only \cmd{/shared} spans both machines. The container home is the GPU
    host's, not the developer VM's. \\[2pt]

    \cmd{<defunct>} processes accumulate &
    Container PID 1 is \cmd{sleep infinity} and never reaps children.
    Harmless; cleared on release. \\\bottomrule
  \end{tabularx}
  \caption{Symptoms seen in practice, with cause and remedy.}
  \label{tab:gpu-pitfalls}
\end{table}

% ============================================================
%  PART II: CORPUS DOCUMENT RETRIEVAL
% ============================================================

\part{Document Retrieval System (corpus)}

\section{Document retrieval on the endpoint: \textit{corpus}}
\label{sec:corpus}

\subsection{What it is}
\textit{corpus} is a retrieval-augmented question answering system over our own
documents, running entirely on American Terawatt hardware and using the
endpoint of Section~\ref{sec:architecture} as its model. Documents are parsed
once, split into passages, embedded, and indexed on local disk. At question
time the relevant passages are retrieved and handed to \cmd{gpt-oss:120b} on
the shared GPU pool, which answers from those passages and cites them. Nothing
leaves the network, and nothing is sent to an external API.

The system deliberately does \emph{not} fine-tune a model on the documents.
Fine-tuning gives no reliable recall of specific facts, no citations, and would
need re-running whenever a document changes. Retrieval at query time gives all
three.

\subsection{Where everything runs}
Three machines are involved, and the split matters: the corpus lives next to
the documents, the model lives on the GPU, and only questions and answers cross
the network.

\begin{figure}[!ht]
  \centering
  \resizebox{\linewidth}{!}{%
  \begin{tikzpicture}[node distance=9mm,
    blk/.style={draw=atBorderStrong, rounded corners=2pt, fill=atBg,
                align=center, inner sep=4pt, font=\scriptsize,
                minimum height=8mm},
    store/.style={blk, fill=atBgSubtle},
    comp/.style={blk, fill=atBgMuted},
    model/.style={blk, fill=atAccent!30},
    zone/.style={draw=atBorderStrong, dashed, rounded corners=4pt, inner sep=6mm},
    flow/.style={{-Stealth[length=2mm]}, draw=atFgMuted, thick},
    net/.style={{-Stealth[length=2mm]}, draw=atFgMuted, thick, dashed},
    lbl/.style={font=\tiny, text=atFgMuted, align=center},
  ]
  % laptop
  \node[comp] (hermes) {Hermes Agent\\{\normalfont (or any MCP client)}};
  \node[zone, fit=(hermes), label={[font=\small\bfseries]above:Your laptop}] (zlap) {};

  % developer VM
  \node[comp, right=30mm of hermes] (server) {corpus server\\{\cmd{:8765}}\\MCP + JSON};
  \node[comp, above right=6mm and 12mm of server] (cli) {{\cmd{./rag} CLI}};
  \node[comp, below right=6mm and 12mm of server] (core)
    {{\textbf{ragkit}}\\parse $\cdot$ chunk\\retrieve $\cdot$ answer};
  \node[store, below=of core] (lance) {LanceDB index\\vectors + BM25};
  \node[store, below=of lance] (md) {{\cmd{md/}} parse cache\\keyed on sha256};
  \node[store, left=9mm of md] (raw) {{\cmd{\char`~/work/raw}}\\dashboard uploads};
  \node[zone, fit=(server)(cli)(core)(lance)(md)(raw),
        label={[font=\small\bfseries]above:dev-eugen {\normalfont(developer VM)}}] (zdev) {};

  % GPU host
  \node[comp, right=30mm of core] (ollama) {Ollama\\{\cmd{:11434}}};
  \node[model, above right=4mm and 10mm of ollama] (chat)
    {{\cmd{gpt-oss:120b}}\\65.4\,GB MXFP4\\131k ctx, tools};
  \node[model, below right=4mm and 10mm of ollama] (embed) {{\cmd{bge-m3}}\\1024-dim};
  \node[zone, fit=(ollama)(chat)(embed),
        label={[font=\small\bfseries]above:GPU host {\normalfont(pool container)}}] (zgpu) {};

  % edges
  \draw[net] (hermes.east) -- node[lbl, above=1mm, midway, text width=24mm]
    {HTTP + token\\office LAN} (server.west);
  \draw[flow] (server.east) -- (core.north west);
  \draw[flow] (cli.south east) -- (core.north);
  \draw[flow] (core) -- (lance);
  \draw[flow] (md) -- (lance);
  \draw[flow] (raw) -- (md);
  \draw[flow] (core.east) -- node[lbl, above=1mm, midway, text width=26mm]
    {HTTP, forwarded\\to the pool} (ollama.west);
  \draw[flow, <->] (ollama.east) -- (chat.west);
  \draw[flow, <->] (ollama.east) -- (embed.west);
  \end{tikzpicture}}
  \caption{Where everything runs. Only questions and answers cross the network;
           the index stays beside the documents.}
  \label{fig:corpus-system}
\end{figure}

\subsection{Ingest: documents in, index out}
Ingest runs once per document, and is cached. Parsing is the slow step, so its
output is stored in \cmd{md/} keyed on the source file's sha256 --- re-running
\cmd{./rag ingest} over a growing folder only touches files whose bytes
changed.

\begin{figure}[!ht]
  \centering
  \begin{tikzpicture}[node distance=6mm,
    blk/.style={draw=atBorderStrong, rounded corners=2pt, fill=atBg,
                align=center, inner sep=4pt, font=\scriptsize,
                minimum height=8mm},
    store/.style={blk, fill=atBgSubtle},
    comp/.style={blk, fill=atBgMuted},
    model/.style={blk, fill=atAccent!30},
    flow/.style={{-Stealth[length=2mm]}, draw=atFgMuted, thick},
    lbl/.style={font=\tiny, text=atFgMuted, text width=26mm},
  ]
  \node[store] (src) {{\cmd{\char`~/work/raw}}\\PDF DOCX XLSX\\CSV MD TXT};
  \node[comp, right=of src] (disc) {discover\\skips venvs,\\\cmd{node\_modules}, \cmd{.git}};
  \node[comp, right=of disc] (parse) {parse\\to markdown};
  \node[comp, right=of parse] (fix) {repair\\ligatures, tidy};
  \node[store, right=of fix] (cache) {{\cmd{md/}}\\cache};

  \node[comp, below=14mm of cache] (chunk) {chunk\\1000 tok, 150 overlap\\heading-aware};
  \node[model, left=of chunk] (emb) {embed\\\cmd{bge-m3}};
  \node[store, left=of emb] (tbl) {LanceDB\\\cmd{chunks} table};
  \node[comp, left=of tbl] (idx) {build indexes\\BM25 always\\ANN $\geq$ 2048 rows};

  \draw[flow] (src) -- (disc);
  \draw[flow] (disc) -- (parse);
  \draw[flow] (parse) -- (fix);
  \draw[flow] (fix) -- (cache);
  \draw[flow] (cache.south) -- (chunk.north);
  \draw[flow] (chunk) -- (emb);
  \draw[flow] (emb) -- (tbl);
  \draw[flow] (tbl) -- (idx);

  \node[lbl, below=1mm of parse] {scans flagged,\\not indexed};
  \node[lbl, below=1mm of chunk, text width=30mm]
    {title + heading path\\prefixed to each chunk};
  \end{tikzpicture}
  \caption{Ingest. Parsing is cached on the source's sha256, so re-running only
           touches documents whose bytes changed.}
  \label{fig:corpus-ingest}
\end{figure}

\subsection{Query: three model calls, not one}
A question triggers an embedding call, then a rerank call, then the answer
call. The database and the model never talk to each other; Python moves text
between them.

\begin{figure}[!ht]
  \centering
  \resizebox{\linewidth}{!}{%
  \begin{tikzpicture}[
    blk/.style={draw=atBorderStrong, rounded corners=2pt, fill=atBg,
                align=center, inner sep=4pt, font=\scriptsize,
                minimum height=8mm},
    comp/.style={blk, fill=atBgMuted},
    model/.style={blk, fill=atAccent!30},
    zone/.style={draw=atBorderStrong, dashed, rounded corners=4pt, inner sep=6mm},
    flow/.style={{-Stealth[length=2mm]}, draw=atFgMuted, thick},
    lbl/.style={font=\tiny, text=atFgMuted, align=center},
  ]
  % retrieval
  \node[blk]   (q)    at (0,0)      {question};
  \node[model] (qe)   at (2.5,0)    {\textbf{1.} embed\\\cmd{bge-m3}};
  \node[comp]  (vec)  at (5.6,1.15) {vector search\\cosine, top 40};
  \node[comp]  (bm)   at (5.6,-1.15){BM25 search\\top 40};
  \node[comp]  (rrf)  at (9.1,0)    {fuse\\reciprocal rank};
  \node[comp]  (prom) at (12.1,0)   {promote a\\dominant\\BM25 hit};
  \node[zone, fit=(vec)(bm),
        label={[font=\tiny, text=atFgMuted]below:LanceDB --- one table, two indexes}] (lz) {};

  % generation
  \node[model] (rr)  at (12.1,-4.3) {\textbf{2.} rerank\\\cmd{gpt-oss:120b}\\listwise, 1 call};
  \node[model] (ans) at (8.0,-4.3)  {\textbf{3.} answer\\\cmd{gpt-oss:120b}};
  \node[blk]   (out) at (4.2,-4.3)  {cited answer\\+ sources};

  \draw[flow] (q) -- (qe);
  \draw[flow] (qe.east) to[out=40,in=180]  (vec.west);
  \draw[flow] (qe.east) to[out=-40,in=180] (bm.west);
  \draw[flow] (vec.east) to[out=0,in=130]  (rrf.north);
  \draw[flow] (bm.east)  to[out=0,in=-130] (rrf.south);
  \draw[flow] (rrf) -- (prom);
  \draw[flow] (prom) -- node[lbl, right=1mm, midway, text width=20mm]{40\\candidates} (rr);
  \draw[flow] (rr) -- node[lbl, above=1mm, midway]{top 8 passages} (ans);
  \draw[flow] (ans) -- (out);

  \node[lbl, text width=40mm] at (5.6,2.5) {embeddings blur exact identifiers};
  \node[lbl, text width=30mm] at (1.5,-1.9) {BM25 catches ``754'', ``IEC 60287''};
  \node[lbl, below=2mm of ans, text width=38mm]
    {{\cmd{--fast}} skips step 2:\\$\approx$7\,s instead of $\approx$20\,s};
  \end{tikzpicture}}
  \caption{A query. Three model calls: embed the question, rerank the
           shortlist, then answer from the survivors.}
  \label{fig:corpus-query}
\end{figure}

\subsection{Model operations}

Four ways to interface with the corpus system:
\begin{itemize}
  \item \cmd{./rag ask} - CLI, no server
  \item Python \cmd{import ragkit} - notebooks, own scripts
  \item MCP stdio - client spawns it
  \item MCP + JSON over HTTP \cmd{:8765} - Hermes and other agents
\end{itemize}

\begin{table}[!ht]
  \centering
  \renewcommand{\arraystretch}{1.25}
  \begin{tabularx}{\linewidth}{@{}l l l X@{}}
    \toprule
    \textbf{Interface} & \textbf{Server?} & \textbf{Auth} & \textbf{Use it for} \\\midrule
    \cmd{./rag ask}              & no      & ---    & everyday use on the developer VM \\\midrule
    \cmd{import ragkit}          & no      & ---    & notebooks, own scripts \\\midrule
    MCP stdio                    & spawned & ---    & Claude Code on that box \\\midrule
    MCP \cmd{/mcp}               & yes     & bearer & Hermes and other agents \\\midrule
    JSON \cmd{/ask}, \cmd{/search} & yes   & bearer & curl, scripts, any language \\\bottomrule
  \end{tabularx}
  \caption{The four interfaces and what each is for.}
  \label{tab:corpus-interfaces}
\end{table}

\subsection{Everyday commands}

\begin{table}[!ht]
  \centering
  \renewcommand{\arraystretch}{1.25}
  \begin{tabularx}{\linewidth}{@{}l X@{}}
    \toprule
    \textbf{Task} & \textbf{Command} \\\midrule
    Index new uploads    & \cmd{./rag ingest} \\\midrule
    Index elsewhere      & \cmd{./rag ingest /shared/eugen/brochures} \\\midrule
    Re-parse everything  & \cmd{./rag ingest --force} \\\midrule
    Rebuild from scratch & \cmd{./rag ingest --reset} \\\midrule
    Ask                  & \cmd{./rag ask "\dots"} \\\midrule
    Passages only        & \cmd{./rag search "\dots"} \\\midrule
    What is indexed      & \cmd{./rag status} \\\midrule
    Start the server     & \cmd{tmux new -s corpus-mcp}, then \cmd{./mcp-http --host 0.0.0.0} \\\midrule
    Client details       & \cmd{./mcp-http --show-client-config} \\\bottomrule
  \end{tabularx}
  \caption{Everyday \textit{corpus} commands on the developer VM.}
  \label{tab:corpus-commands}
\end{table}

\section{Known limits}

\begin{itemize}\itemsep1pt
  \item \textbf{Scanned PDFs index nothing.} Ingest flags them rather than
        silently adding empty chunks. A vision model (\cmd{qwen2.5vl:32b}) is
        now on the pool and is the obvious OCR path if needed.
  \item \textbf{Damage already in the source is not repaired.} The eCIGRE
        brochure list contains titles that lost their ligatures before reaching
        the spreadsheet (``Model specication''), so a search for
        ``specification'' will not match them. Ligature repair runs only for
        PDFs, where the document's own text supplies a trustworthy vocabulary.
  \item \textbf{One writer.} LanceDB is embedded: concurrent readers are fine,
        two simultaneous ingests are not.
  \item \textbf{Rerank latency.} The 120B spends reasoning tokens before any
        visible output, so a reranked query costs $\approx$20\,s.
  \item \textbf{LAN exposure.} With \cmd{--host 0.0.0.0} the bearer token is
        the only thing between the office network and every indexed document.
\end{itemize}

% ============================================================
%  PART III: APPENDIX - TIKZ TRANSMISSION LINE DIAGRAM
% ============================================================

\appendix

\chapter{TikZ Transmission Line Diagram}

\section{AC Transmission Line Model}

The following TikZ diagram represents the transmission line model from image9.pdf:

\begin{figure}[!ht]
  \centering
  \begin{circuitikz}[american voltages, scale=1.2]
    % Source box
    \draw (0,0) rectangle (1,1);
    
    % Label above source (P_S, Q_S)
    \draw[->,violet, thick] (1,1.3) -- (3,1.3) node[midway, above] {$P_S, Q_S$};
    
    % Line from source to point S
    \draw (1,0.5) -- (2,0.5);
    
    % Point S (red dot)
    \fill[red] (2,0.5) circle (2pt);
    \node[right] at (2,0.6) {$S$};
    
    % I_S arrow (red, going right)
    \draw[->,red, thick] (2,0.5) -- (2.8,0.5) node[midway, above] {$I_S$};
    
    % U_S arrow (blue, going down)
    \draw[->,blue, thick] (2,0.5) -- (2,-0.5) node[midway, right] {$U_S$};
    
    % Switch/open
    \draw[thick] (3,0.5) -- (3.3,0.5);
    \draw[thick] (3.3,0.5) -- (3.5,0.6) -- (3.5,0.4) -- (3.8,0.5);
    
    % Cable (thick black line) with brace
    \draw[ultra thick] (3.8,0.5) -- (6,0.5);
    \draw[decorate,decoration={brace,amplitude=6pt}, above] 
        (3.8,0.5) -- (6,0.5) node[midway,above=8pt] {$L_{\text{cable}}$};
    
    % Point R (red dot)
    \fill[red] (6,0.5) circle (2pt);
    \node[right] at (6,0.6) {$R$};
    
    % I_R arrow (red, going right)
    \draw[->,red, thick] (6,0.5) -- (6.8,0.5) node[midway, above] {$I_R$};
    
    % U_R arrow (blue, going down)
    \draw[->,blue, thick] (6,0.5) -- (6,-0.5) node[midway, right] {$U_R$};
    
    % Label above point R (P_R, Q_R)
    \draw[->,violet, thick] (6,1.3) -- (8,1.3) node[midway, above] {$P_R, Q_R$};
    
    % Switch/open
    \draw[thick] (6.8,0.5) -- (7.2,0.5);
    \draw[thick] (7.2,0.5) -- (7.5,0.6) -- (7.5,0.4) -- (7.8,0.5);
    
    % Load Z_L
    \draw (7.8,0.5) to[R, l^=$Z_L$] (7.8,-1) node[ground]{};
    
    % Ground from cable
    \draw (6,-0.5) -- (6,-1) node[ground]{};
    \draw (2,-0.5) -- (2,-1) node[ground]{};
  \end{circuitikz}
  \caption{TikZ representation of the transmission line model from image9.pdf.}
  \label{fig:transmission-line-tikz}
\end{figure}

\section{Diagram Description}

The diagram shows:
\begin{itemize}
  \item Source box on the left with power flow labels ($P_S, Q_S$)
  \item Point $S$ (sending end) marked with red dot
  \item Current $I_S$ (red arrow) and voltage $U_S$ (blue arrow)
  \item Cable segment $L_{\text{cable}}$ with thick line and brace
  \item Point $R$ (receiving end) marked with red dot
  \item Current $I_R$ (red arrow) and voltage $U_R$ (blue arrow)
  \item Power flow labels ($P_R, Q_R$) at receiving end
  \item Load $Z_L$ (resistor) to ground
\end{itemize}

\chapter{Command Reference}

\section{GPU Pool Commands}

\begin{lstlisting}
# --- pool -----------------------------------------------------------
gpu status
gpu claim 1
gpu shell                       # interactive
gpu shell /path/to/script.sh    # one command
gpu release

# --- inside the container -------------------------------------------
nvidia-smi
ps -eo pid,etime,cmd | grep "[b]in/ollama serve"
tail -f /shared/$USER/ollama-serve.log

# --- model ----------------------------------------------------------
ollama list
ollama pull  gpt-oss:120b
ollama run   gpt-oss:120b --verbose "prompt"
ollama stop  gpt-oss:120b       # unload from VRAM, keep on disk

# --- endpoint -------------------------------------------------------
curl http://localhost:11435/api/tags
curl http://localhost:11435/api/generate \
  -d '{"model":"gpt-oss:120b","prompt":"...","stream":false}'
\end{lstlisting}

\section{Corpus Commands}

\begin{lstlisting}
# --- index ------------------------------------------------------------
./rag ingest                            # new uploads in ~/work/raw
./rag ingest /shared/eugen/brochures    # a folder elsewhere
./rag ingest --force                    # re-parse everything
./rag ingest --reset                    # rebuild from scratch
./rag status                            # what is indexed

# --- ask --------------------------------------------------------------
./rag ask    "..."                      # cited answer
./rag ask    "..." --fast               # skip the rerank call
./rag search "..."                      # passages only

# --- server -----------------------------------------------------------
tmux new -s corpus-mcp
./mcp-http --host 0.0.0.0               # MCP + JSON on :8765, bearer token
./mcp-http --show-client-config
\end{lstlisting}

\end{document}
